% ch47_straightening.tex — Volume IV Chapter 11

\chapter{Straightening and Unstraightening}

\begin{overviewbox}
The \term{straightening} and \term{unstraightening} constructions of Joyal
and Lurie make precise the relationship previewed in Chapter~46: Cartesian
fibrations over a simplicial set $\mathcal{S}$ are equivalent, in a
homotopy-coherent sense, to functors $\mathcal{S}^{\mathrm{op}} \to
\mathcal{C}\mathrm{at}_\infty$.

The full theorem is one of the heaviest technical results in $\infty$-category
theory — Lurie's \emph{Higher Topos Theory} devotes Chapter~2 to it,
roughly 100 pages of dense combinatorial argument. We follow the chapter's
proof-depth policy: state the theorem clearly, work the low-dimensional
cases explicitly, sketch the technical strategy, and refer to HTT §2.2
for the full apparatus. The reader who has digested this chapter is then
prepared to consult HTT directly.

The chapter introduces:
\begin{itemize}
  \item \term{Marked simplicial sets} $\mathbf{sSet}^+$: simplicial sets
        together with a marked subset $M \subseteq X_1$ containing
        degeneracies. The marked edges are the ``would-be Cartesian
        morphisms.''
  \item The \term{Cartesian model structure} on $\mathbf{sSet}^+\!/\mathcal{S}$
        whose fibrant objects are Cartesian fibrations over $\mathcal{S}$
        (with their Cartesian edges as the marking).
  \item The \term{straightening functor} $\mathrm{St}_\mathcal{S} :
        \mathbf{sSet}^+\!/\mathcal{S} \to \mathrm{Fun}(\mathfrak{C}[\mathcal{S}]^{\mathrm{op}},
        \mathbf{sSet}^+)$, where $\mathfrak{C}[\mathcal{S}]$ is the
        simplicial category rigidification of $\mathcal{S}$.
  \item The \term{unstraightening} $\mathrm{Un}_\mathcal{S}$ as right adjoint.
\end{itemize}

Theorem~\ref{thm:straightening}: $\mathrm{St}_\mathcal{S} \dashv \mathrm{Un}_\mathcal{S}$
is a Quillen equivalence between the Cartesian model structure on
$\mathbf{sSet}^+\!/\mathcal{S}$ and the projective model structure on
$\mathrm{Fun}(\mathfrak{C}[\mathcal{S}]^{\mathrm{op}}, \mathbf{sSet}^+)$.
On homotopy categories, this gives the equivalence
\begin{equation*}
  \mathrm{Cart}(\mathcal{S}) \;\simeq\; \mathrm{Fun}(\mathcal{S}^{\mathrm{op}}, \mathcal{C}\mathrm{at}_\infty),
\end{equation*}
the precise statement of ``Cartesian fibrations $=$ presheaves of
$\infty$-categories.''
\end{overviewbox}


% ═══════════════════════════════════════════════════════════════════════════
\section{The Straightening Equivalence}
\label{sec:straightening-statement}
% ═══════════════════════════════════════════════════════════════════════════

\subsection{Formalizing}

\begin{definition}[Marked simplicial set]
\label{def:marked-sset}
A \term{marked simplicial set} is a pair $(X, M)$ where $X \in \mathbf{sSet}$
and $M \subseteq X_1$ is a subset containing all degenerate $1$-simplices.
A morphism $(X, M) \to (Y, N)$ is a simplicial-set map $X \to Y$ sending
$M$ into $N$. The category is $\mathbf{sSet}^+$.
\end{definition}

\begin{howtosay}
$\mathbf{sSet}^+$ \sayit{``marked simplicial sets''}; $(X, M)$ has $M$ as
the ``marked edges'' or ``would-be equivalences/Cartesian morphisms.'' For
a quasi-category $\mathcal{C}$, the natural marking is its equivalences;
for a Cartesian fibration $p : \mathcal{X} \to \mathcal{S}$, the natural
marking on $\mathcal{X}$ is its Cartesian morphisms.
\end{howtosay}

\begin{definition}[Rigidification $\mathfrak{C}$]
\label{def:rigidification}
The \term{rigidification} functor $\mathfrak{C} : \mathbf{sSet} \to
\mathbf{sCat}$ sends a simplicial set $\mathcal{S}$ to a simplicial category
whose objects are the vertices of $\mathcal{S}$ and whose mapping
simplicial sets $\mathfrak{C}[\mathcal{S}](x, y)$ are specific Kan
complexes built from the simplices of $\mathcal{S}$ with appropriate
end-conditions.

For $\mathcal{S} = \mathrm{N}(\mathcal{D})$ the nerve of an ordinary
category, $\mathfrak{C}[\mathcal{S}]$ is the trivial simplicial enrichment
of $\mathcal{D}$ (mapping sets are discrete simplicial sets).
\end{definition}

\begin{theorem}[Straightening--Unstraightening Equivalence]
\label{thm:straightening}
For every simplicial set $\mathcal{S}$, there is a Quillen adjunction
\begin{equation*}
  \mathrm{St}_\mathcal{S} : \mathbf{sSet}^+ \!/ \mathcal{S} \;\;\rightleftarrows\;\;
  \mathrm{Fun}(\mathfrak{C}[\mathcal{S}]^{\mathrm{op}}, \mathbf{sSet}^+) : \mathrm{Un}_\mathcal{S}
\end{equation*}
between the Cartesian model structure on $\mathbf{sSet}^+\!/\mathcal{S}$
and the projective model structure on $\mathrm{Fun}(\mathfrak{C}[\mathcal{S}]^{\mathrm{op}},
\mathbf{sSet}^+)$. The adjunction is a Quillen equivalence.

On homotopy categories, after restricting to fibrant objects, this gives
\begin{equation*}
  \mathrm{Cart}(\mathcal{S}) \;\simeq\; \mathrm{Fun}(\mathcal{S}^{\mathrm{op}}, \mathcal{C}\mathrm{at}_\infty).
\end{equation*}
\end{theorem}

\begin{proof}[Proof strategy]
The full proof spans HTT §2.2 (about 100 pages of dense combinatorics).
The strategy:

\emph{Step 1: Construct $\mathrm{St}_\mathcal{S}$ explicitly.} For
$(\mathcal{X}, M) \to \mathcal{S}$ a marked simplicial set over
$\mathcal{S}$, define $\mathrm{St}_\mathcal{S}(\mathcal{X}, M)(s)$ as a
specific simplicial set built from the marked $n$-simplices over chains
starting at $s$.

\emph{Step 2: Verify the adjunction.} The right adjoint $\mathrm{Un}_\mathcal{S}$
has an explicit formula as a certain limit; the unit and counit are
checked combinatorially.

\emph{Step 3: Prove Quillen.} Use the cofibrant generation of both model
structures and check that $\mathrm{St}_\mathcal{S}$ sends generating
(trivial) cofibrations to (trivial) cofibrations.

\emph{Step 4: Prove Quillen equivalence.} The hardest step. Verify that
unit and counit on cofibrant–fibrant objects are weak equivalences.
This uses the rigidification's good behavior and a detailed analysis of
how marked simplicial sets present Cartesian fibrations.

We will work through the low-dimensional cases in
\S\ref{sec:straightening-low-dim} below. \qed
\end{proof}


% ═══════════════════════════════════════════════════════════════════════════
\section{Low-Dimensional Warm-Up}
\label{sec:straightening-low-dim}
% ═══════════════════════════════════════════════════════════════════════════

\subsection{Formally: $\mathcal{S} = \Delta^0$}

\begin{example_thm}[Straightening over a point]
\label{ex:straighten-point}
When $\mathcal{S} = \Delta^0$, $\mathfrak{C}[\Delta^0]$ is the trivial
simplicial category with one object and one morphism. Then:
\begin{itemize}
  \item $\mathbf{sSet}^+\!/\Delta^0 = \mathbf{sSet}^+$.
  \item Cartesian fibrations over $\Delta^0$ are simply quasi-categories
        (with their natural marking by equivalences).
  \item $\mathrm{Fun}(\mathfrak{C}[\Delta^0]^{\mathrm{op}}, \mathbf{sSet}^+) =
        \mathbf{sSet}^+$.
\end{itemize}
The straightening is the identity on the underlying $\mathbf{sSet}$.
The equivalence trivially holds.
\end{example_thm}

\subsection{Formally: $\mathcal{S} = \Delta^1$}

\begin{example_thm}[Straightening over the walking arrow]
\label{ex:straighten-arrow}
When $\mathcal{S} = \Delta^1$, $\mathfrak{C}[\Delta^1]$ is a simplicial
category with two objects $0, 1$ and one non-trivial morphism $0 \to 1$
(plus identities). The mapping simplicial set $\mathfrak{C}[\Delta^1](0, 1)$
is $\Delta^0$.

A Cartesian fibration $p : \mathcal{X} \to \Delta^1$ has fibres $\mathcal{X}_0$
and $\mathcal{X}_1$, and the Cartesian condition gives a transport functor
$\phi : \mathcal{X}_1 \to \mathcal{X}_0$ (over the edge $0 \to 1$ in
$\Delta^1$). The data is captured by:
\begin{itemize}
  \item Two quasi-categories: $\mathcal{X}_0, \mathcal{X}_1$.
  \item A functor: $\phi : \mathcal{X}_1 \to \mathcal{X}_0$.
\end{itemize}

A functor $\mathfrak{C}[\Delta^1]^{\mathrm{op}} \to \mathbf{sSet}^+$ is
also exactly the data of: $F(1), F(0)$ (two marked simplicial sets) and
a morphism $F(0 \to 1)^{\mathrm{op}} = \phi : F(1) \to F(0)$.

So $\mathrm{St}_{\Delta^1}$ assigns $(\mathcal{X}_0, \mathcal{X}_1, \phi)$
to $\mathcal{X} \to \Delta^1$, which is the same data on both sides. The
straightening is essentially the identity. \qed
\end{example_thm}

\subsection{Formally: general low-dimensional pattern}

\begin{remark}[Why low-dimensional examples are easy]
For $\mathcal{S}$ a $1$-dimensional simplicial set (so $\mathfrak{C}[\mathcal{S}]$
is a trivial simplicial category), the straightening reduces to ordinary
$1$-categorical Grothendieck unstraightening. Coherence comes for free
because all mapping simplicial sets in $\mathfrak{C}[\mathcal{S}]$ are
$\Delta^0$.

The non-trivial content of $\mathrm{St}_\mathcal{S}$ appears for
$\mathcal{S}$ of dimension $\geq 2$ (where $\mathfrak{C}[\mathcal{S}]$ has
non-trivial mapping simplicial sets) and especially for $\mathcal{S}$
a general quasi-category (where the rigidification injects
infinite-dimensional combinatorics from $2$-simplices and above).
\end{remark}


% ═══════════════════════════════════════════════════════════════════════════
\section{The Cartesian Model Structure}
\label{sec:cartesian-model}
% ═══════════════════════════════════════════════════════════════════════════

\subsection{Formalizing}

\begin{theorem}[Cartesian model structure on $\mathbf{sSet}^+\!/\mathcal{S}$]
\label{thm:cartesian-model}
The category $\mathbf{sSet}^+\!/\mathcal{S}$ admits a model structure —
the \term{Cartesian model structure} — characterized by:
\begin{itemize}
  \item Cofibrations: monomorphisms (of underlying simplicial sets).
  \item Fibrant objects: Cartesian fibrations $(\mathcal{X}, M_{\mathrm{Cart}}) \to \mathcal{S}$,
        marked by their Cartesian edges.
  \item Weak equivalences: maps inducing equivalences on the corresponding
        $\infty$-categories of straightened functors.
\end{itemize}
\end{theorem}

\begin{remark}[Two markings on a quasi-category]
A quasi-category $\mathcal{C}$ has two natural markings:
\begin{itemize}
  \item Maximal marking $\mathcal{C}^\sharp$: every $1$-simplex marked.
  \item Minimal (equivalence) marking $\mathcal{C}^\flat$: only
        equivalences marked.
\end{itemize}
The maximal marking corresponds to viewing $\mathcal{C}$ as a Kan complex
(every edge is "an equivalence" by fiat); the minimal marking is the
``honest'' one for $\infty$-category theory. The Cartesian model structure
sees them differently.
\end{remark}

\subsection{Reorganizing: presheaves as a model}

\begin{theorem}[Equivalence with presheaves of $\infty$-categories]
\label{thm:equiv-presheaves}
On homotopy categories, the straightening induces an equivalence
\begin{equation*}
  \mathrm{Ho}(\mathbf{sSet}^+\!/\mathcal{S})_{\mathrm{Cart}} \;\simeq\;
  \mathrm{Ho}\bigl(\mathrm{Fun}(\mathfrak{C}[\mathcal{S}]^{\mathrm{op}}, \mathbf{sSet}^+)\bigr)_{\mathrm{proj}}.
\end{equation*}
Both sides are models for the homotopy category of presheaves of
$\infty$-categories on $\mathcal{S}$. Cartesian fibrations are an
``encoded form'' of such presheaves.
\end{theorem}


% ═══════════════════════════════════════════════════════════════════════════
\section{Computational Use}
\label{sec:straightening-computational}
% ═══════════════════════════════════════════════════════════════════════════

\subsection{Formally}

\begin{example_thm}[Pulling back classifying functors]
\label{ex:pullback-classifying}
Given a functor $F : \mathcal{S}^{\mathrm{op}} \to \mathcal{C}\mathrm{at}_\infty$,
the corresponding Cartesian fibration $\mathcal{X}_F \to \mathcal{S}$
satisfies: for every $s \in \mathcal{S}_0$, the fibre $(\mathcal{X}_F)_s
\simeq F(s)$ in $\mathcal{C}\mathrm{at}_\infty$. The unstraightening
$\mathrm{Un}_\mathcal{S}$ implements this passage.

Pullback: for $f : \mathcal{T} \to \mathcal{S}$, the composition
$\mathrm{Un}_\mathcal{T}(F \circ f^{\mathrm{op}}) = \mathcal{X}_F \times_\mathcal{S}
\mathcal{T}$.
\end{example_thm}

\begin{example_thm}[Representable presheaves and slices]
\label{ex:representable-slice}
For an object $x$ of a quasi-category $\mathcal{C}$, the slice fibration
$\mathcal{C}_{/x} \to \mathcal{C}$ is a right fibration whose classifying
functor under straightening is the representable presheaf
\begin{equation*}
  \mathrm{Map}_\mathcal{C}(-, x) : \mathcal{C}^{\mathrm{op}} \to \mathcal{S}.
\end{equation*}
Right fibrations land in the subcategory of \emph{$\infty$-groupoid-valued}
functors $\mathcal{C}^{\mathrm{op}} \to \mathcal{S}$. The straightening
gives, in particular, a precise sense in which ``every right fibration
is a presheaf of spaces.''
\end{example_thm}

\subsection{Reorganizing: where this goes}

\begin{remark}[Phase-2 milestone]
With straightening in place, Chapter~48 will state and prove the
$\infty$-categorical Yoneda lemma in the form
\begin{equation*}
  \mathcal{C} \;\hookrightarrow\; \mathcal{P}(\mathcal{C}) := \mathrm{Fun}(\mathcal{C}^{\mathrm{op}}, \mathcal{S}),
\end{equation*}
a fully faithful embedding into presheaves of spaces. This is the
$\infty$-categorical avatar of the classical Yoneda embedding (Chapter~9,
15, 20, 28), and it organizes the rest of Volume~IV.
\end{remark}


% ═══════════════════════════════════════════════════════════════════════════
\section{Exercises}
\label{sec:straightening-exercises}
% ═══════════════════════════════════════════════════════════════════════════

\begin{exercisesbox}
\begin{qa}
\qaitem{What does the straightening theorem say in one sentence?}{Cartesian fibrations over $\mathcal{S}$ are equivalent (as $\infty$-categories) to functors $\mathcal{S}^{\mathrm{op}} \to \mathcal{C}\mathrm{at}_\infty$.}
\qaitem{What is a marked simplicial set?}{A pair $(X, M)$ with $X \in \mathbf{sSet}$ and $M \subseteq X_1$ containing degeneracies. The marking is the would-be Cartesian/equivalence morphisms.}
\qaitem{Why mark simplicial sets at all?}{Because Cartesian fibrations and their morphisms are best handled when the Cartesian edges are explicitly tracked. The Cartesian model structure on $\mathbf{sSet}^+$ uses this marking.}
\qaitem{What is the rigidification $\mathfrak{C}[\mathcal{S}]$?}{A simplicial category with the same objects as $\mathcal{S}$, whose mapping simplicial sets are explicit Kan complexes built from $\mathcal{S}$'s simplices. Reduces to discrete enrichment when $\mathcal{S} = \mathrm{N}(\mathcal{D})$.}
\qaitem{State the straightening adjunction.}{$\mathrm{St}_\mathcal{S} : \mathbf{sSet}^+\!/\mathcal{S} \rightleftarrows \mathrm{Fun}(\mathfrak{C}[\mathcal{S}]^{\mathrm{op}}, \mathbf{sSet}^+) : \mathrm{Un}_\mathcal{S}$, Quillen equivalence.}
\qaitem{What do the fibres of a Cartesian fibration become under straightening?}{The value of the corresponding functor $\mathcal{S}^{\mathrm{op}} \to \mathcal{C}\mathrm{at}_\infty$ at each vertex.}
\qaitem{What's the straightening over $\Delta^0$?}{Trivially the identity: Cartesian fibrations over a point are quasi-categories, and functors from the one-object category to $\mathbf{sSet}^+$ are just $\mathbf{sSet}^+$.}
\qaitem{What's the straightening over $\Delta^1$?}{The data of a Cartesian fibration $p : \mathcal{X} \to \Delta^1$ — two fibres $\mathcal{X}_0, \mathcal{X}_1$ and a functor $\mathcal{X}_1 \to \mathcal{X}_0$ — equals the data of a functor $\mathfrak{C}[\Delta^1]^{\mathrm{op}} \to \mathbf{sSet}^+$. Both are: ``two marked simplicial sets and a morphism between them.''}
\qaitem{Why is straightening hard for higher-dimensional $\mathcal{S}$?}{Because $\mathfrak{C}[\mathcal{S}]$ has non-trivial mapping simplicial sets when $\mathcal{S}$ has $2$-simplices and above. The combinatorics of comparing fibrations and functors then becomes substantial.}
\qaitem{What are the two natural markings on a quasi-category $\mathcal{C}$?}{Maximal ($\mathcal{C}^\sharp$): all edges marked. Minimal ($\mathcal{C}^\flat$): only equivalences marked. The two correspond to different categorical pictures (Kan-complex view vs.\ honest $\infty$-category view).}
\qaitem{What is the Cartesian model structure on $\mathbf{sSet}^+\!/\mathcal{S}$?}{Cofibrations: monomorphisms; fibrant objects: Cartesian fibrations with their Cartesian markings; weak equivalences: fibrewise equivalences plus base.}
\qaitem{How does the equivalence read at the level of homotopy categories?}{$\mathrm{Cart}(\mathcal{S}) \simeq \mathrm{Fun}(\mathcal{S}^{\mathrm{op}}, \mathcal{C}\mathrm{at}_\infty)$ — the homotopy category of Cartesian fibrations over $\mathcal{S}$ is the same as the homotopy category of $\mathcal{S}^{\mathrm{op}}$-indexed diagrams of $\infty$-categories.}
\qaitem{What about coCartesian fibrations?}{The dual theorem: coCartesian fibrations over $\mathcal{S}$ are equivalent to functors $\mathcal{S} \to \mathcal{C}\mathrm{at}_\infty$ (covariant). Symmetric under $\mathcal{S} \leftrightarrow \mathcal{S}^{\mathrm{op}}$.}
\qaitem{And right fibrations?}{Right fibrations correspond to functors $\mathcal{S}^{\mathrm{op}} \to \mathcal{S}$ (into the $\infty$-category of spaces). Because right fibrations have all-Cartesian morphisms and Kan-complex fibres.}
\qaitem{What does the slice $\mathcal{C}_{/x}$ straighten to?}{The representable presheaf $\mathrm{Map}_\mathcal{C}(-, x) : \mathcal{C}^{\mathrm{op}} \to \mathcal{S}$.}
\qaitem{How is straightening used in practice?}{To convert between fibration-form and functor-form of structured $\infty$-categorical data. Many constructions are easier to describe as fibrations (universal example: limits, colimits) but easier to manipulate as functors (universal example: adjoint functor existence).}
\qaitem{What's the universal Cartesian fibration?}{$\mathcal{U} \to \mathcal{C}\mathrm{at}_\infty^{\mathrm{op}}$, whose fibre over $\mathcal{D}$ is $\mathcal{D}$. Straightens to $\mathrm{id} : \mathcal{C}\mathrm{at}_\infty^{\mathrm{op}} \to \mathcal{C}\mathrm{at}_\infty$.}
\qaitem{Why is the full proof of straightening so long?}{Because verifying that $\mathrm{St}_\mathcal{S}$ is a Quillen equivalence requires checking unit and counit on fibrant–cofibrant objects involve intricate combinatorial decomposition of marked simplicial sets and their relationship to the rigidification.}
\qaitem{Where in the literature is the full proof?}{Lurie, \emph{Higher Topos Theory}, §2.2 (about 100 pages). Riehl–Verity, \emph{Elements of $\infty$-Category Theory}, gives an alternative model-independent treatment.}
\qaitem{Can one avoid straightening?}{Yes — Riehl–Verity develop $\infty$-category theory model-independently, using comma constructions and never explicitly invoking the straightening equivalence. Practical for many purposes, but the equivalence itself is a fundamental structural fact.}
\qaitem{Why does this chapter sketch rather than prove?}{Per the Volume~IV proof-depth policy (selective full proofs, otherwise sketches). The straightening proof itself fills a textbook chapter and would dilute the bridge-to-literature mission.}
\qaitem{What chapter uses straightening crucially?}{Chapter~48 ($\infty$-Yoneda) --- the embedding $\mathcal{C} \to \mathcal{P}(\mathcal{C})$ is via straightening of the slice fibrations.}
\qaitem{What is the homotopy-categorical statement to take away?}{$\mathrm{Cart}(\mathcal{S}) \simeq \mathrm{Fun}(\mathcal{S}^{\mathrm{op}}, \mathcal{C}\mathrm{at}_\infty)$. This is the $\infty$-categorical Grothendieck construction.}
\end{qa}
\end{exercisesbox}


% ═══════════════════════════════════════════════════════════════════════════
\section*{Chapter Summary}
\addcontentsline{toc}{section}{Chapter Summary}
% ═══════════════════════════════════════════════════════════════════════════

\begin{summarybox}
\textbf{Marked simplicial sets.}\quad $\mathbf{sSet}^+ = $ pairs $(X, M)$
with $M \subseteq X_1$ containing degeneracies; the marking encodes
Cartesian/equivalence edges.

\textbf{Rigidification.}\quad $\mathfrak{C} : \mathbf{sSet} \to
\mathbf{sCat}$; sends a simplicial set to a simplicial category with
specific Kan-complex mapping spaces. Trivial on nerves of $1$-categories.

\textbf{Straightening--unstraightening.}\quad For every $\mathcal{S}$,
$\mathrm{St}_\mathcal{S} \dashv \mathrm{Un}_\mathcal{S}$ is a Quillen
equivalence between $(\mathbf{sSet}^+\!/\mathcal{S})_{\mathrm{Cart}}$
and $\mathrm{Fun}(\mathfrak{C}[\mathcal{S}]^{\mathrm{op}}, \mathbf{sSet}^+)_{\mathrm{proj}}$.

\textbf{Homotopy-category statement.}\quad
\begin{equation*}
  \mathrm{Cart}(\mathcal{S}) \;\simeq\; \mathrm{Fun}(\mathcal{S}^{\mathrm{op}}, \mathcal{C}\mathrm{at}_\infty).
\end{equation*}
``Cartesian fibrations $=$ presheaves of $\infty$-categories.''

\textbf{Low-dimensional reduction.}\quad For $\mathcal{S}$ a nerve of a
$1$-category, straightening reduces to ordinary Grothendieck construction.
The full $\infty$-categorical content appears for $\mathcal{S}$ with
non-trivial mapping simplicial sets in its rigidification.

\textbf{Use.}\quad Chapter~48 uses straightening to construct the
$\infty$-Yoneda embedding $\mathcal{C} \hookrightarrow \mathcal{P}(\mathcal{C})$.
\end{summarybox}


% ═══════════════════════════════════════════════════════════════════════════
\section*{Formal Restatement}
\addcontentsline{toc}{section}{Formal Restatement}
% ═══════════════════════════════════════════════════════════════════════════

\begin{formalrestatement}
\textbf{Marked simplicial sets.}\quad $\mathbf{sSet}^+ = $ category of pairs
$(X, M)$, $M \supseteq \{\text{degenerate } 1\text{-simplices of } X\}$.

\medskip
\textbf{Rigidification.}\quad $\mathfrak{C} : \mathbf{sSet} \to \mathbf{sCat}$,
the left adjoint to the simplicial nerve $\mathrm{N}^{\mathrm{coh}}$.

\medskip
\textbf{Cartesian model structure on $\mathbf{sSet}^+\!/\mathcal{S}$.}\quad
\begin{itemize}
  \item Cofibrations: monomorphisms.
  \item Fibrant: Cartesian fibrations $(\mathcal{X}, M_{\mathrm{Cart}}) \to \mathcal{S}$.
  \item Weak equivalences: maps inducing equivalences on classifying functors.
\end{itemize}

\medskip
\textbf{Straightening--unstraightening (Theorem~\ref{thm:straightening}).}\quad
$\mathrm{St}_\mathcal{S} \dashv \mathrm{Un}_\mathcal{S}$, Quillen equivalence:
\begin{equation*}
  \mathbf{sSet}^+\!/\mathcal{S} \;\simeq_{\mathrm{Q}}\; \mathrm{Fun}\bigl(\mathfrak{C}[\mathcal{S}]^{\mathrm{op}}, \mathbf{sSet}^+\bigr).
\end{equation*}

\medskip
\textbf{Homotopy-category statement.}\quad
$\mathrm{Cart}(\mathcal{S}) \simeq \mathrm{Fun}(\mathcal{S}^{\mathrm{op}}, \mathcal{C}\mathrm{at}_\infty)$;
$\mathrm{coCart}(\mathcal{S}) \simeq \mathrm{Fun}(\mathcal{S}, \mathcal{C}\mathrm{at}_\infty)$;
$\mathrm{Right}(\mathcal{S}) \simeq \mathrm{Fun}(\mathcal{S}^{\mathrm{op}}, \mathcal{S})$;
$\mathrm{Left}(\mathcal{S}) \simeq \mathrm{Fun}(\mathcal{S}, \mathcal{S})$.

\medskip
\textbf{Slice as representable.}\quad $\mathcal{C}_{/x} \to \mathcal{C}$
corresponds under straightening to $\mathrm{Map}_\mathcal{C}(-, x) :
\mathcal{C}^{\mathrm{op}} \to \mathcal{S}$.
\end{formalrestatement}
